\section{電子計算機のイロハ}
\label{D1_lesson02}

\subsection{授業内容}
\subsubsection{科目の中でのこのコマの位置づけ}
\begin{tabular}[t]{ccccccc}
    記述統計[0] & 確率[0] & モデル[0] & 推測・推定[0] & 意思決定[0] & 実践[0]
\end{tabular}
\subsubsection{概要}
昨今はスマートフォン，タブレットなどのデバイスが行き渡っているが，専門的な計算をするためにはやはりパソコンを使うことになる。この講義は次回からのPCを使った演習に対する準備的内容であり，必ずしも心理統計に直結する内容ではない。とはいえ，統計計算の歴史は計算機発展の歴史でもあり，実際にPCのもつ基本的な知識・概念・技能がないと以降のPCを使った演習にも差し障りが出るので，十分な準備をしてから実践に進む必要がある。本講は心理統計の内容に直接関係するものではないが，今後の実習・演習に必要な知識の事前準備をする。
\subsubsection{コマ主題細目(箇条書き)}
\begin{description}
    \item[コンピュータの基礎] 21世紀の私たちは身の回りをコンピュータに囲まれて生きている。それは携帯電話の形をしていたり，タブレットの形をしていたり，ノートパソコン，デスクトップパソコンなどの形をしているだろう。他にも富嶽のような大型計算機もあるが，これらの計算機はいずれも入力，出力，記憶，演算，制御，記憶という5つの機能を持った複合体である。コンピュータの使い方として，入出力装置の使い方はもちろん，一時記憶装置，外部記憶装置の区別などを理解する。またGUIによる直感的な操作は，その背後機械的な操作があり，ユーザには操作した感覚というイリュージョンを見せているだけであることに注意する。
    \begin{flushright}
        $\to$この件に関しては，心理学の意識論とも関係の深い\citet{トールノーレットランダーシュ2002}も参照
    \end{flushright}
    \item[コンピュータの歴史] コンピュータ開発の歴史は小型化の歴史でもある。大きくこの流れを見るとともに，OSの歴史についても概観し，現在のコンピュータ環境，OSとアプリケーションの違い，フリーソフトウェアとシェアウェア(プロプライエタリ)の違いなどにも言及する。
    \begin{flushright}
        $\to$フリーソフトウェアの考え方については\citet{ストールマン2003}を，Macの歴史については\citet{アイザックソン2012}，日本のネット社会の歴史については\citet{杉本2020}など，参考になる読み物は色々と挙げられる。
    \end{flushright}
     
    \item[情報の単位] コンピュータが様々な形をとりながらも情報をやり取りできているのは，テキスト，画像，音声，動画などあらゆるものを「情報」という単位に落とし込んで一元的に管理しているからである。このことに加えて，ここでは情報の単位を理解する。また文字コードによるデータの欠落の問題についても注意する。そのための半角英数文字によるファイル名の付け方が重要であることを指摘する。
     
    \item[ファイルの場所]あらゆる媒体が情報という一元化された単位で扱われるのは便利なことではあるが，一見してその違いがわからないのは不便である。そこでファイル名には拡張子と呼ばれるファイルの種類を区別する情報を付与する。また，ファイルがたくさん出来上がって管理できないのも困るので，いくつかのファイルをまとめるフォルダ/ディレクトリという単位でまとめて置くことを理解する。またこれらのファイルがどこにあるかをしっかり理解する必要がある。特にクラウド/ローカルの違いについて注意しておく。
    
    \item[再現性のある研究実践へ]今後，パソコンを使って文章のみならず，図表も作成することになる。統計的分析結果の出力も，プログラムによって行うことになる。コンピュータにおいてプログラムが実行されるとはどういうことか(適切なファイルへのアクセス)を理解し，またプログラムを書いて実行するときの注意点について理解する。特にエラーとその報告の仕方については，恐怖心に駆られることなく対応するべきである。そして研究の実践とその成果が「動作」ではなく「指示書」として記録に残ることは，研究の再現性を高めるためにも重要な営みであることを理解する。Gitによる文書のバージョン管理や，OSFなど情報の共有など，最新の研究実践方法についても言及する。

\end{description}
\subsubsection{キーワード}
\begin{itemize}
\item コンピュータの5大機能
\item ビット，バイト
\item 文字コード
\item フォルダと拡張子
\item アルゴリズム
\end{itemize}
\subsection{授業情報}
\paragraph{コマの展開方法}
講義
\paragraph{標準シラバスにおける位置づけ}
科目番号5；心理学統計法；2.統計に関する基礎的な知識；C エクセル，R，SPSS入門(１) 統計分析のための言語の手ほどき　プログラムの基本的操作
\subsubsection{予習・復習課題}
\paragraph{予習}大学生活ではPCを使うことがほぼ確実に必要であり，今後の社会においてもその能力は間違いなく生きてくるので，なるべく早い段階でPCを所持し，ブラインドタッチなど基本的な技術を習得しておいて欲しい。
\paragraph{復習}自分のPCのユーザ名，ファイル名，拡張子の付け方などを確認し，半角英数になっていないか確認しよう。今後ファイルを作ることがあれば，半角英数文字で名前をつけるよう癖づけておいたほうが良い。これらは未然にエラーを防ぐことにもなる。
